\chapter{Conclusion}
Retailers hold CCTV footage that could answer everyday operational questions but is rarely used for that purpose, and small and mid-sized stores often cannot afford specialised analytics. RetailPulse addresses this with a CPU-friendly, non-biometric, multi-agent design that reuses existing cameras. Research on detection, tracking, and re-identification, and a feasibility triage, shaped the scope: a single-store, three-camera prototype with a conservative two-camera handoff in place of full re-identification.

At the mid-report stage the requirements and architecture are defined. The web application with manager and employee roles, the POS analytics service, secure session handling, report export, and the camera-handoff MVP are implemented. The floor-distance CLI now reuses the trained employee/customer detector and ByteTrack, exports provisional floor positions/distances, and rejects unsupported alignment. Seventeen focused checks and four existing handoff regression tests passed, and sampled actual-video exports were verified. Dashboard vision metrics still rely on prototype data; the CLI exports approximate floor-distance observations and their quality states. The main risks are CPU throughput, classification accuracy, and cross-camera matching uncertainty. FYP-2 will integrate the vision agents through the LangGraph orchestrator, evaluate them on labelled store footage, and harden the system for the final demonstration.
